\documentclass{exam}
\usepackage[left=0.25in,right=0.25in,top=1in,bottom=1in]{geometry}
\usepackage{amsmath}
\usepackage{graphicx}
\usepackage{nicematrix}
\usepackage{extpfeil}
\usepackage{amsthm}
\usepackage{gensymb}
\usepackage{amsfonts}
\usepackage{tikz}
\usepackage{pgfplots}
\usepackage{array}
\usepackage{hyperref}
\usepackage{nameref}
\usepackage[x11names]{xcolor}
\usepackage[most]{tcolorbox}

\pgfplotsset{compat=1.18}
\printanswers

\newcommand{\ZZ}{\mathbb Z}
\newcommand{\QQ}{\mathbb Q}
\newcommand{\NN}{\mathbb N}
\newcommand{\CC}{\mathbb C}
\newcommand{\RR}{\mathbb R}
\newcommand{\braces}[1]{\ensuremath{\left\{#1 \right\}}}
\newcommand{\Span}[1]{\ensuremath\text{Span}\braces{#1}}
\newcommand{\ee}{\mathbf{e}}

\hypersetup{colorlinks=true, linktoc=section, linkcolor=blue}

\pagestyle{headandfoot}
\firstpageheadrule
\runningheadrule
\firstpageheader{Prof. Shen \\ Differential Equations}{Formula Sheet}{Jeevan Shah}
\runningheader{Differential Equations \\ Formula Sheet}{}{Shah}
\firstpagefooter{}{\thepage}{}
\runningfooter{ }{\thepage}{ }

\printanswers


\begin{document}
    \underline{Matrix Exponential:} Given a square matrix $A$ the matrix exponential, $e^{tA}$, can be calculated by diagonalizing $A$ into $A=PDP^{-1}$ such that if $A$ has eigenvalues $\lambda_1, \lambda_2, \ldots, \lambda_n$ with corresponding eigenvectors $u_1, u_2, \ldots, u_n$ then 
    \[
        e^{tA} = Pe^{tD}P^{-1} = \begin{pmatrix}
            u_1 & u_2 & \cdots & u_n
        \end{pmatrix}  
        \begin{pmatrix}
            e^{t\lambda_1} & 0 & \cdots & 0 \\
            0 & e^{t\lambda_2} & \cdots & 0 \\
            \vdots & & \ddots &  \vdots \\
            0 & \cdots & 0 & e^{\lambda_n}
        \end{pmatrix}
        \begin{pmatrix}
            u_1 & u_2 & \cdots & u_n
        \end{pmatrix}^{-1}
    \]
    If $A$ has a repeated eigenvalue $\lambda$, with corresponding eigenvector $v_1$, then the generalized eigenvector can be found by solving $(A - \lambda I)^{k}v_2 = v_1$ for some $k$. Standardly we have $(A - \lambda I)v_2 = v_1$ where $v_2$ is the generalized eigenvector. As well, the Jordan form will change to 
    \[
        J = \begin{pmatrix}
            \lambda & 1 \\
            0 & \lambda
        \end{pmatrix} 
        \Rightarrow e^{tJ} = \begin{pmatrix}
            e^{\lambda} & te^{\lambda} \\
            0 & e^{\lambda}
        \end{pmatrix}
    \]
    We also have that for 
    \[
        \mathbf{x}'(t) = A\mathbf{x}(t), \quad \mathbf{x}(0) = \mathbf{x}_0
    \]
    the general solution is 
    \[
        \mathbf{x}(t) = e^{tA}\mathbf{x}_0
    \]
    with $e^{tA}$ the matrix exponential.

    \underline{Duhamels Formula:} For the non-autonomous equation 
    \[
        \mathbf{x}'(t) = A\mathbf{x}(t) + \mathbf{g}(t)
    \]
    \textit{Duhamels Formula} provides the general solution: 
    \[
        \mathbf{x}(t) = e^{(t - t_0)A}\mathbf{x}_0 + \int_{t_0}^t e^{(t-s)A}\mathbf{g}(s)\,\mathrm{d}s
    \]
    for the initial value $\mathbf{x}(t_0)= \mathbf{x}_0$.

    \underline{Picard Iteration:} For an IVP with 
    \[
        x'(t) = v(x(t), t), \quad x(t_0) = x_0 
    \]
    define 
    \[
        X_0 = x_0,\,
        X_1 = x_0 + \int_{t_0}^{t} v(x_0(s), s)\, \mathrm{d}s, \,\ldots,\,
        X_n = x_0 + \int_{t_0}^{t} v(x_{n-1}(s), s)\, \mathrm{d}s.
    \]
    The solution to the IVP is then $\lim_{n\to\,\infty} X_n$ 

    \underline{Supremum Distance:} For $f, g$ continuous functions on $[a, b]$ define 
    \[
        d(f, g) = \max_{t \in [a, b]} |f(t) - g(t)|.
    \]
    Tip: Square the distance to remove the absolute value bars

    \underline{Second Order ODEs:} For an equation in the form 
    \[
        ax''(t) + bx'(t) + cx(t) = f(t) 
    \]
    start by solving the homogenous equation where $f(t) = 0$ and $a, b, c \in \RR$ , with the characteristic equation $ar^2 + br + c = 0$. 

    \begin{minipage}{0.45\linewidth} 
        \begin{itemize}
            \item If $r_1 \neq r_2$ then $x_h = C_{1}e^{r_{1}t}+ C_{2}e^{r_{2}t}$
        \end{itemize}
    \end{minipage}
    \hfill
    \begin{minipage}{0.45\linewidth}
        \begin{itemize}
            \item If $r_1 = r_2$ then $x_h = (C_1 + C_2t)e^{r_{1}t}$
        \end{itemize}
    \end{minipage}

    Then solve for the particular solution by guessing:

    \begin{itemize}
        \item $g(x) = e^{\alpha t} \Rightarrow x_p = Ce^{\alpha t}$
        \item $g(x)$ is an $n$-degree polynomial $\Rightarrow x_p = A_{n}t^{n} + \cdots + A_{1}t + A_0$ (can ommit constant term based on context)
        \item $g(x) = A\cos bt$ or $A\sin bt \Rightarrow x_p = C \cos bt + D \sin bt$.
    \end{itemize}        
    If the $x_p$ guess is already accounted for in the homogenous solution, multiply by $t$ until it is unique. The final solution is then $x = x_p + x_h$.

    \newpage 

    \underline{Cauchy-Euler Equations}: For an equation in the form 
    \[
        t^2 x'' + at x' + bx = g(t) 
    \]
    guess $x_1 = t^{\alpha}$ for the homogenous equation so that $x_1' = \alpha t^{\alpha -1}$ and $x_1'' = \alpha(\alpha - 1)t^{\alpha - 2}$. It follows that 
    \[
        \alpha(\alpha-1) + a\alpha + b = 0    
    \]

    \begin{minipage}{0.45\linewidth}
        \begin{itemize}
            \item If $\alpha_1 \neq \alpha_2$ then $x_h = C_1 t^{\alpha_{1}} + C_2 t^{\alpha_{2}}$
        \end{itemize}    
    \end{minipage}
    \hfill
    \begin{minipage}{0.45\linewidth}
        \begin{itemize}
            \item If $\alpha_1 = \alpha_2$ then $x_h = C_1 t^{\alpha_{1}} + C_2 t^{\alpha_{1}} \ln t$
        \end{itemize}    
    \end{minipage}

    Then find the particular solution by guessing or variation of parameters. The final solution is then $x = x_1 + x_p$.

    \underline{Variation of Parameters:} If $x_1(t)$ and $x_2(t)$ are solutions to the homogenous equation then a particular solution is 
    \[
        x_p = \int_{t_{0}}^{t} = \frac{\det(N(t, s))}{\det(M(s))}f(s)\,\mathrm{d}s
    \]
    for 
    \[
        N(t, s) = \begin{pmatrix}
            x_1(s) & x_2(s) \\
            x_1(t) & x_2(t)
        \end{pmatrix}    
        \quad\text{and}\quad
        M(s) = \begin{pmatrix}
            x_1(s) & x_2(s) \\
            x_1'(s) & x_2'(s) 
        \end{pmatrix}
    \]

    \underline{Solution Space Dim:} A $n$th order differential equation has a solution space with order $n$. Thus, any solution set with more than $n$ solutions, must contain at least one linearly dependent solution.


    \underline{Inner Product Spaces:} A function $f$ defined on $[a, b]$ is in $L^{2}[a, b]$, the inner product space of all square integrable functions, if, and only if, 
    \[
        \int_{a}^{b} |f(x)|^2\,\mathrm{d}x < \infty 
    \]

    \underline{Approximation in Inner Product Spaces:} To find coefficents $c_0, c_1, c_2, \ldots, c_k$ to approximation some function $f$ in $L^{2}[a, b]$ given a set of basis functions $\braces{\phi_1, \phi_2, \ldots \phi_k}$, we must minimize 
    \[
        \left\|f - (c_0 + c_1\phi_1 + c_2\phi_2 + \cdots + c_k\phi_k)\right\|^2 = \int_{a}^{b} (f - c_0 - c_1\phi_1 - c_2\phi_2 - \cdots - c_k\phi_k)^2\,\mathrm{d}x.
    \]
    We can find the smallest such $c_k$ by computing the projection: 
    \[
        c_k = \frac{\langle f, \phi_k\rangle}{\langle \phi_k, \phi_k \rangle} 
    \]
    where the inner product $\langle f, g \rangle = \int_{a}^{b} f(x)g(x)\,\mathrm{d}x$ provided the basis functions are all orthogonal (i.e. $\langle f, g \rangle = \int_{a}^{b} f(x)g(x)\,\mathrm{d}x = 0$ for all basis functions).
    
    \underline{Good to know formulas:}

    \begin{minipage}{0.45\linewidth}
        \[
            e^{t} = \sum_{k=0}^{\infty} \frac{t^{k}}{k\mathbf{!}} = 1 + t + \frac{t^2}{2\mathbf{!}} + \frac{t^2}{3\mathbf{!}} + \cdots 
        \]
    \end{minipage}
    \hfill
    \begin{minipage}{0.45\linewidth}
        \[
            te^{t} = \sum_{k=0}^{\infty} \frac{t^{k+1}}{k\mathbf{!}} = t + t^2 + \frac{t^3}{2\mathbf{!}} + \frac{t^3}{3\mathbf{!}} + \cdots 
        \]
    \end{minipage}
    \begin{minipage}{0.45\linewidth}
        \[
            \sin t = \sum_{k=0}^{\infty} \frac{(-1)^{k}}{(2k + 1)\mathbf{!}}t^{2k + 1} = t - \frac{t^3}{3\mathbf{!}} + \frac{t^5}{5\mathbf{!}} + \cdots 
        \]
    \end{minipage}
    \hfill
    \begin{minipage}{0.45\linewidth}
        \[
            \cos t = \sum_{k=0}^{\infty} \frac{(-1)^{k}}{(2k +)\mathbf{!}}t^{2k} = 1 - \frac{t^3}{3\mathbf{!}} + \frac{t^4}{4\mathbf{!}} + \cdots 
        \]
    \end{minipage}
    \begin{minipage}{0.45\linewidth}
        \[
            \sin 2t = 2 \sin t \cos t
        \]
    \end{minipage}
    \hfill
    \begin{minipage}{0.45\linewidth}
        \[
            \cos 2t = \cos^2 t - \sin^2 t
        \]
    \end{minipage}
    \begin{minipage}{0.45\linewidth}
        \[
            \sin(\alpha \pm \beta) = \sin \alpha \cos \beta \pm \cos\alpha \sin\beta
        \]
    \end{minipage}
    \hfill
    \begin{minipage}{0.45\linewidth}
        \[
            \cos(\alpha \pm \beta) = \cos\alpha\cos\beta + \sin\alpha\sin\beta
        \]
    \end{minipage}
    \[
        \frac{1}{1-t} = \sum_{t=0}^{\infty} t^n = 1 + t + t^2 + \cdots 
    \]
    \begin{minipage}{0.45\linewidth}
        \[
            \cos^{2}(x) = \frac{1+\cos(2x)}{2}
        \]
    \end{minipage}
    \hfill
    \begin{minipage}{0.45\linewidth}
        \[
            \sin^{2}(x) = \frac{1-\cos(2x)}{2}
        \]
    \end{minipage}
    \begin{minipage}{0.45\linewidth}
        \[
            A = \begin{pmatrix}
                a & b \\
                c & d
            \end{pmatrix}
            \Rightarrow A^{-1} = \frac{1}{ad - bc}\begin{pmatrix}
                d & -b \\
                -c & a
            \end{pmatrix}
        \]
    \end{minipage}
    \hfill
    \begin{minipage}{0.45\linewidth}
        \[
            A = \begin{pmatrix}
                a & b \\
                c & d
            \end{pmatrix} 
            \Rightarrow \lambda^2 - (a + d)\lambda + (ad - bc) = 0
        \]
    \end{minipage}
\end{document}